%%
%% part2-book-components/ch05-cover.tex
%%
%% Chapter 5: Cover and Back Cover — Full-color front cover
%% and static two-block back cover with author bio.
%%

\chapter{جلد و پشت جلد}
\label{chap:cover}

\section{چرا این موضوع مهم است؟}
\label{sec:cover-why}

جلد کتاب، اولین چیزی است که خواننده می‌بیند و
نقش مهمی در جذب مخاطب دارد. در کتاب‌های فنی و
علمی، جلد حرفه‌ای نشان‌دهنده‌ی کیفیت محتواست و
اعتماد خواننده را جلب می‌کند.

\lr{MatinBook} یک سیستم جلد کامل ارائه می‌دهد که
شامل جلد جلو تمام‌رنگی و جلد عقب با دو بلوک
ایستاتیک است. هدف این فصل، آموزش نحوه‌ی استفاده
و شخصی‌سازی این سیستم است.

\mnote{جلد \lr{MatinBook} با \lr{TikZ} ساخته می‌شود و به حداقل دو بار کامپایل نیاز دارد.}

\section{جلد جلو}
\label{sec:cover-front}

جلد جلو با دستور \cmd{makecover} ساخته می‌شود.
این دستور چهار آرگومان می‌گیرد:

\begin{latin}
\begin{minted}{latex}
\makecover
    {|\pc{عنوان کتاب}|}
    {|\pc{زیرعنوان کتاب}|}
    {|\pc{نام نویسنده}|}
    {|\pc{مهر ۱۴۰۵}|}
\end{minted}
\end{latin}

\begin{enumerate}
    \item \textbf{عنوان کتاب:} عنوان اصلی.
    \item \textbf{زیرعنوان:} توضیح کوتاه یا شعار.
    \item \textbf{نام نویسنده:} نویسنده یا نویسندگان.
    \item \textbf{تاریخ یا نسخه:} تاریخ یا شماره‌ی نسخه.
\end{enumerate}

\mnote{اگر زیرعنوان خالی باشد، از \lr{\{\}} استفاده کنید: \lr{\textbackslash makecover\{عنوان\}\{\}\{نویسنده\}\{تاریخ\}}.}

\subsection{عناصر بصری جلد جلو}
\label{sec:cover-front-elements}

جلد جلو \lr{MatinBook} از عناصر زیر تشکیل شده
است:

\begin{itemize}
    \item \textbf{پس‌زمینه‌ی سرمه‌ای:} رنگ اصلی جلد
    که با \lr{matinnavy} تعریف شده است.
    \item \textbf{گوه‌ی قرمز:} یک مثلث قرمز در
    گوشه‌ی بالا-چپ.
    \item \textbf{بیضی طلایی:} یک بیضی طلایی
    در گوشه‌ی پایین-راست.
    \item \textbf{شبکه‌ی مختصات:} یک شبکه‌ی
    کم‌کنتراست برای بافت.
    \item \textbf{نمونه‌کد C++:} یک نمونه‌ی کد
    با فریم طلایی.
    \item \textbf{معادلات ریاضی:} چهار معادله‌ی
    معروف.
    \item \textbf{فلش‌ها و نقاط:} عناصر تزئینی.
\end{itemize}

\mnote{همه‌ی این عناصر در فایل \lr{mb-theme-cover.sty} تعریف شده‌اند. برای شخصی‌سازی، این فایل را ویرایش کنید.}

\subsection{نمای ساده‌شده‌ی جلد}
\label{sec:cover-front-mockup}

شکل \cref{fig:cover-mockup} نمای ساده‌شده‌ای
از جلد جلو را نشان می‌دهد:

\begin{figure}[h]
\centering
\begin{latin}
\begin{tikzpicture}[scale=0.5]
  % Background
  \fill[matinnavy] (0,0) rectangle (10,14);
  % Crimson wedge
  \fill[matincrimson] (0,14) -- (5.5,14) -- (3,11.5) -- cycle;
  % White frame
  \draw[white, opacity=0.6, line width=0.6pt]
      (0.4,0.4) rectangle (9.6,13.6);
  % Gold ellipse
  \draw[covergold, opacity=0.4, line width=1pt, rotate=-18]
      (7,3) ellipse (3.5 and 1.8);
  % Title area
  \node[white, font=\Large\bfseries] at (7,9) {\rl{عنوان کتاب}};
  \draw[covergold, line width=1pt] (6,8) -- (9,8);
  \node[white, font=\small] at (7.5,7.5) {\rl{زیرعنوان}};
  % Author
  \node[white, font=\bfseries] at (7,4.5) {\rl{نام نویسنده}};
  % MatinBook label
  \node[white, font=\footnotesize, anchor=west] at (1,0.8) {MatinBook};
  \node[white, font=\footnotesize, anchor=east] at (9,0.8) {github.com/...};
\end{tikzpicture}
\end{latin}
\caption{نمای ساده‌شده‌ی جلد جلو}
\label{fig:cover-mockup}
\end{figure}

\mnote{در \lr{tikzpicture}، متن فارسی باید داخل \lr{\textbackslash rl\{...\}} قرار گیرد.}

\subsection{شخصی‌سازی جلد جلو}
\label{sec:cover-front-custom}

برای شخصی‌سازی جلد جلو، سه راه دارید:

\textbf{راه ۱ — تغییر آرگومان‌ها:} فقط محتوای
\cmd{makecover} را تغییر دهید.

\textbf{راه ۲ — تغییر رنگ‌ها:} در
\file{mb-theme-cover.sty}، رنگ‌های
\lr{coverprimary}، \lr{coveraccent}، \lr{covergold}
و \lr{coverpaper} را تغییر دهید.

\textbf{راه ۳ — تغییر طرح:} فایل
\file{mb-theme-cover.sty} را ویرایش کنید و
عناصر جدید اضافه یا حذف کنید.

\mnote{برای تغییرات ساده، راه ۱ و ۲ کافی است. برای طراحی کاملاً جدید، راه ۳ را انتخاب کنید.}

\section{جلد عقب}
\label{sec:cover-back}

جلد عقب با دستور \cmd{makebackcover} ساخته
می‌شود. این دستور یک آرگومان می‌گیرد:

\begin{latin}
\begin{minted}{latex}
\makebackcover{|\pc{توضیحات پشت جلد...}|}
\end{minted}
\end{latin}

\mnote{دستور \lr{\textbackslash makebackcover} باید در انتهای سند، پیش از \lr{\textbackslash end\{document\}} قرار گیرد.}

\subsection{ساختار دو بلوکی}
\label{sec:cover-back-two-blocks}

جلد عقب \lr{MatinBook} از دو بلوک ایستاتیک
تشکیل شده است:

\begin{enumerate}
    \item \textbf{بلوک اول — درباره این کتاب:}
    با عنوان \lr{\textbackslash coverabout} و متن
    توضیحات.
    \item \textbf{بلوک دوم — درباره این نویسنده:}
    با عنوان \lr{\textbackslash coverauthorbio} و
    متن زندگی‌نامه.
\end{enumerate}

فاصله‌ی بین دو بلوک، دقیقاً \pnum{۰.۷}
سانتی‌متر است.

\mnote{هر دو بلوک در یک \lr{minipage} با ارتفاع ثابت قرار می‌گیرند تا متن طولانی باعث سرریز به صفحه‌ی دوم نشود.}

\subsection{افزودن زندگی‌نامه نویسنده}
\label{sec:cover-back-authorbio}

برای افزودن زندگی‌نامه‌ی نویسنده، از دستور
\cmd{authorbio} استفاده کنید:

\begin{latin}
\begin{minted}{latex}
\authorbio{%
    |\pc{نام نویسنده، پژوهشگر حوزه‌ی الگوریتم و ریاضیات است.}|%
    |\pc{او سال‌ها در زمینه‌ی آموزش برنامه‌نویسی فعالیت کرده است.}|%
}

\makebackcover{%
    |\pc{توضیحات پشت جلد...}|%
}
\end{minted}
\end{latin}

\mnote{اگر \lr{\textbackslash authorbio} را فراخوانی نکنید، بلوک دوم به‌کلی حذف می‌شود و فقط بلوک کتاب نمایش داده می‌شود.}

\subsection{شخصی‌سازی جلد عقب}
\label{sec:cover-back-custom}

برای شخصی‌سازی جلد عقب:

\begin{itemize}
    \item \textbf{تغییر عنوان‌ها:} با
    \cmd{renewcommand} می‌توانید
    \lr{\textbackslash coverabout} و
    \lr{\textbackslash coverauthorbio} را تغییر دهید.
    \item \textbf{تغییر رنگ‌ها:} در
    \file{mb-theme-cover.sty}، رنگ‌های
    \lr{coveraccent} و \lr{coverprimary} را
    تغییر دهید.
    \item \textbf{تغییر متن:} فقط آرگومان‌های
    \cmd{makebackcover} و \cmd{authorbio} را
    تغییر دهید.
\end{itemize}

\mnote{عنوان‌های \lr{\textbackslash coverabout} و \lr{\textbackslash coverauthorbio} در فایل \lr{fa-IR.sty} تعریف شده‌اند.}

\section{اطلاعات مخزن}
\label{sec:cover-repository}

جلد \lr{MatinBook} آدرس مخزن پروژه را نمایش
می‌دهد. برای تغییر آن:

\begin{latin}
\begin{minted}{latex}
\renewcommand{\repository}{github.com/username/repo}
\end{minted}
\end{latin}

\mnote{مقدار پیش‌فرض \lr{github.com/matinmahmoudi-cs/matinbook} است.}

\section{اطلاعات نسخه}
\label{sec:cover-version}

جلد \lr{MatinBook} نسخه‌ی کلاس را نمایش می‌دهد.
این مقدار در \file{matinbook.cls} تعریف شده است:

\begin{latin}
\begin{minted}{latex}
\def\matinbookversion{1.2}
\def\matinbookdate{2026/10/08}
\end{minted}
\end{latin}

\mnote{نسخه‌ی \lr{MatinBook} در صفحه‌ی شناسنامه و روی جلد نمایش داده می‌شود.}

\section{رنگ‌های جلد}
\label{sec:cover-colors}

جلد \lr{MatinBook} از رنگ‌های زیر استفاده می‌کند
که در \file{mb-theme-cover.sty} تعریف شده‌اند:

\begin{matintable}{رنگ‌های جلد \lr{MatinBook}}{tab:cover-colors}
\begin{tblr}{
    width = \textwidth,
    colspec = {l X[l] X[c]},
    hlines, vlines,
    colsep = 8pt,
    rowsep = 4pt,
    row{1} = {font=\bfseries},
}
نام & نقش & مقدار \lr{RGB} \\
\lr{coverprimary} & پس‌زمینه‌ی جلد جلو & \lr{(30, 42, 58)} \\
\lr{coveraccent} & گوه و عناصر تزئینی & \lr{(166, 25, 46)} \\
\lr{covergold} & بیضی و متن طلایی & \lr{(201, 162, 39)} \\
\lr{coverpaper} & پس‌زمینه‌ی جلد عقب & \lr{(242, 238, 227)} \\
\lr{covergray} & متن‌های ثانویه & \lr{(145, 151, 157)} \\
\end{tblr}
\end{matintable}

\mnote{این رنگ‌ها با پالت \lr{CMYK} کتاب هماهنگ هستند تا جلد با محتوای داخلی یکپارچه باشد.}

\section{کامپایل جلد}
\label{sec:cover-compile}

جلد \lr{MatinBook} با \lr{TikZ} ساخته می‌شود و
از \lr{remember picture, overlay} استفاده می‌کند.
این یعنی حداقل \textbf{دو بار کامپایل} لازم است:

\begin{latin}
\begin{minted}{bash}
# |\pc{مرحله‌ی اول: نوشتن موقعیت}| \lr{TikZ}
xelatex -shell-escape document.tex

# |\pc{مرحله‌ی دوم: استفاده از موقعیت}|
xelatex -shell-escape document.tex
\end{minted}
\end{latin}

\mnote{اگر جلد ناقص یا جابه‌جا ظاهر شد، یک بار دیگر کامپایل کنید.}

\section{خطاهای رایج}
\label{sec:cover-mistakes}

\subsection{قرار دادن \cmd{makebackcover} در ابتدای سند}
\label{sec:cover-mistake-position}

اگر \cmd{makebackcover} را در ابتدای سند قرار
دهید، جلد عقب روی صفحه‌ی دوم ظاهر می‌شود و
بقیه‌ی کتاب به‌هم می‌ریزد.

\textbf{راه‌حل:} همیشه \cmd{makebackcover} را در
انتهای سند قرار دهید.

\subsection{فراموش کردن \cmd{authorbio}}
\label{sec:cover-mistake-authorbio}

اگر \cmd{authorbio} را فراموش کنید، بلوک «درباره
این نویسنده» حذف می‌شود و جلد عقب فقط یک بلوک
خواهد داشت.

\textbf{راه‌حل:} اگر می‌خواهید زندگی‌نامه‌ی نویسنده
نمایش داده شود، \cmd{authorbio} را پیش از
\cmd{makebackcover} قرار دهید.

\subsection{سرریز متن جلد عقب به صفحه‌ی دوم}
\label{sec:cover-mistake-overflow}

اگر متن جلد عقب بسیار طولانی باشد، ممکن است
به صفحه‌ی دوم سرریز کند. \lr{MatinBook} این
مشکل را با \lr{minipage} با ارتفاع ثابت حل
کرده است، ولی اگر متن بسیار طولانی باشد،
ممکن است همچنان سرریز رخ دهد.

\textbf{راه‌حل:} متن را کوتاه‌تر کنید یا از
\cmd{small} برای کاهش اندازه‌ی فونت استفاده کنید.

\section{تمرین‌ها}
\label{sec:cover-exercises}

\begin{exercise}
    یک جلد جلو با عنوان، زیرعنوان، نام نویسنده
    و تاریخ بسازید. آیا همه‌ی عناصر به‌درستی
    نمایش داده می‌شوند؟
    \label{exc:cover-front}
\end{exercise}

\begin{exercise}
    یک جلد عقب با دو بلوک بسازید و از
    \cmd{authorbio} استفاده کنید.
    \label{exc:cover-back}
\end{exercise}

\begin{exercise}
    رنگ‌های جلد را تغییر دهید و کتاب را
    دوباره کامپایل کنید. آیا رنگ‌ها اعمال
    می‌شوند؟
    \label{exc:cover-colors}
\end{exercise}

\begin{exercise}
    یک جلد عقب با متن بسیار طولانی بسازید.
    آیا به صفحه‌ی دوم سرریز می‌کند؟
    \label{exc:cover-overflow}
\end{exercise}

\section{خلاصه}
\label{sec:cover-summary}

در این فصل، سیستم جلد \lr{MatinBook} را بررسی
کردیم:

\begin{itemize}
    \item جلد جلو با \cmd{makecover} و چهار
    آرگومان ساخته می‌شود.

    \item جلد جلو شامل پس‌زمینه‌ی سرمه‌ای، گوه‌ی
    قرمز، بیضی طلایی و عناصر تزئینی است.

    \item جلد عقب با \cmd{makebackcover} و یک
    آرگومان ساخته می‌شود.

    \item جلد عقب دو بلوک ایستاتیک دارد:
    «درباره این کتاب» و «درباره این نویسنده».

    \item زندگی‌نامه‌ی نویسنده با \cmd{authorbio}
    اضافه می‌شود.

    \item جلد با \lr{TikZ} ساخته می‌شود و به
    حداقل دو بار کامپایل نیاز دارد.

    \item خطاهای رایج شامل قرار دادن
    \cmd{makebackcover} در ابتدای سند، فراموش
    کردن \cmd{authorbio} و سرریز متن است.
\end{itemize}

در فصل بعدی (\cref{chap:chapters})، با فصل‌ها
و بخش‌ها آشنا می‌شوید و یاد می‌گیرید که چگونه
ساختار کتاب خود را سازمان‌دهی کنید.